\section{AMC 10 10B 2003}
        \begin{enumerate}
        	\item (\textit{AMC 10 10B 2003}) Which of the following is the same as


 
 \[\frac{2-4+6-8+10-12+14}{3-6+9-12+15-18+21}?\]


 $\textbf{(A)} -1 \qquad\textbf{(B)} -\frac{2}{3} \qquad\textbf{(C)} \frac{2}{3} \qquad\textbf{(D)} 1 \qquad\textbf{(E)} \frac{14}{3}$


 
 \textbf{Answer: }C



\textbf{Solution 1}


 \[\frac{2-4+6-8+10-12+14}{3-6+9-12+15-18+21}=\frac{2+(-4+6)+(-8+10)+(-12+14)}{3+(-6+9)+(-12+15)+(-18+21)}=\frac{2\cdot4}{3\cdot4}=\boxed{\text {(C) }\ \frac{2}{3}}\]

 



\textbf{Solution 2}

Each term in the numerator is $\text {(C) } \frac{2}{3}$ of a corresponding term in the denominator, so this must be the quotient.

 


	\item (\textit{AMC 10 10B 2003}) Al gets the disease algebritis and must take one green pill and one pink pill each day for two weeks. A green pill costs $\$$1$ more than a pink pill, and Al's pills cost a total of $\$$546$ for the two weeks. How much does one green pill cost?


 $\textbf{(A) }\$$7 \qquad\textbf{(B) }\$ $14 \qquad\textbf{(C) }\$$19\qquad\textbf{(D) }\$ $20\qquad\textbf{(E) }\$$39$


 
 \textbf{Answer: }D



\textbf{Solution}

Because there are $14$ days in two weeks, Al spends $546/14 = 39$ dollars per day for the cost of a green pill and a pink pill. If the green pill costs $x$ dollars and the pink pill $x-1$ dollars, the sum of the two costs $2x-1$ should equal $39$ dollars. Then the cost of the green pill $x$ is $\boxed{\textbf{(D) }\textdollar  20}$.

 


	\item (\textit{AMC 10 10B 2003}) The sum of $5$ consecutive even integers is $4$ less than the sum of the first $8$ consecutive odd counting numbers. What is the smallest of the even integers?


 $\textbf{(A) } 6 \qquad\textbf{(B) } 8 \qquad\textbf{(C) } 10 \qquad\textbf{(D) } 12 \qquad\textbf{(E) } 14$


 
 \textbf{Answer: }B



\textbf{Solution}

It is a well-known fact that the sum of the first $n$ odd numbers is $n^2$.  This makes the sum of the first $8$ odd numbers equal to $64$.

 Let $n$ be equal to the smallest of the $5$ even integers.  Then $(n+2)$ is the next highest, $(n+4)$ even higher, and so on.

 This sets up the equation $n+(n+2)+(n+4)+(n+6)+(n+8)+4=64$

 Now we solve:

 
 \[5n+24=64\]
 \[5n=40\]
 \[n=8\]

 Thus, the smallest integer is $\boxed{\textbf{(B) } 8}$.

 


	\item (\textit{AMC 10 10B 2003}) Rose fills each of the rectangular and square regions of her rectangular flower bed with a different type of flower. The lengths, in feet, of the rectangular regions in her flower bed are as shown in the figure. She plants one flower per square foot in each region. Asters cost $\$$1$ each, begonias $\$$1.50$ each, cannas $\$$2$ each, dahlias $\$$2.50$ each, and Easter lilies $\$$3$ each. What is the least possible cost, in dollars, for her garden?


 
                    \begin{figure}[H]
                        \centering
                        \includegraphics[width=0.5\textwidth]{Problem/AMC_10_10B_2003_4_Problem_Figure_1.png}
                    \end{figure}
                    


 $\textbf{(A) } 108 \qquad\textbf{(B) } 115 \qquad\textbf{(C) } 132 \qquad\textbf{(D) } 144 \qquad\textbf{(E) } 156$


 
 \textbf{Answer: }A



\textbf{Solution}

The areas of the five regions from greatest to least are $21,20,15,6$ and $4$.

 If we want to minimize the cost, we want to maximize the area of the cheapest flower and minimize the area of the most expensive flower. Doing this, the cost is $1\cdot21+1.50\cdot20+2\cdot15+2.50\cdot6+3\cdot4$, which simplifies to $\$$108$.Therefore the answer is $\boxed{\textbf{(A) } 108}$.

 


	\item (\textit{AMC 10 10B 2003}) Moe uses a mower to cut his rectangular $90$-foot by $150$-foot lawn. The swath he cuts is $28$ inches wide, but he overlaps each cut by $4$ inches to make sure that no grass is missed. He walks at the rate of $5000$ feet per hour while pushing the mower. Which of the following is closest to the number of hours it will take Moe to mow the lawn?


 $\textbf{(A) } 0.75 \qquad\textbf{(B) } 0.8 \qquad\textbf{(C) } 1.35 \qquad\textbf{(D) } 1.5 \qquad\textbf{(E) } 3$


 
 \textbf{Answer: }C



\textbf{Solution}

Since the swath Moe actually mows is $24$ inches, or $2$ feet wide, he mows $10000$ square feet in one hour. His lawn has an area of $13500$, so it will take Moe $1.35$ hours to finish mowing the lawn. Thus the answer is $\boxed{\textbf{(C) } 1.35}$.

 


 



\textbf{Solution 2 (very easy to understand)}

Let's assume that the swath moves back and forth; parallel to the $90$ feet side. Thus, the length of one strip is $90$ feet. Now we need to find out how many strips there are. In reality, the swath Moe mows is $24$ inches wide, which can be easily translated into $2$ feet. $\frac{150}{2}$ is the number of strips Moe needs to mow, which is equal to $75$. Therefore, the total number of feet Moe mows is $75\times90$. Since Moe's mowing rate is $5000$ feet per hour, $\frac{75\times90}{5000}$ is the number of hours it takes him to do his job. Using basic calculations, we compute the answer. $\boxed{\textbf{(C) } 1.35}$.

 


	\item (\textit{AMC 10 10B 2003}) Many television screens are rectangles that are measured by the length of their diagonals. The ratio of the horizontal length to the height in a standard television screen is $4 : 3$. The horizontal length of a "$27$-inch" television screen is closest, in inches, to which of the following?


 
                    \begin{figure}[H]
                        \centering
                        \includegraphics[width=0.5\textwidth]{Problem/AMC_10_10B_2003_6_Problem_Figure_1.png}
                    \end{figure}
                    
$\textbf{(A) } 20 \qquad\textbf{(B) } 20.5 \qquad\textbf{(C) } 21 \qquad\textbf{(D) } 21.5 \qquad\textbf{(E) } 22$


 
 \textbf{Answer: }D



\textbf{Solution 1}

If you divide the television screen into two right triangles, the legs are in the ratio of $4 : 3$, and we can let one leg be $4x$ and the other be $3x$. Then we can use the Pythagorean Theorem.

 
 \begin{align*}(4x)^2+(3x)^2&=27^2\\ 16x^2+9x^2&=729\\ 25x^2&=729\\ x^2&=\frac{729}{25}\\ x&=\frac{27}{5}\\ x&=5.4\end{align*}


 The horizontal length is $5.4\times4=21.6$, which is closest to $\boxed{\textbf{(D) \ } 21.5}$.

 



\textbf{Solution 2}

One can realize that the diagonal, vertical, and horizontal lengths all make up a $3,4,5$ triangle. Therefore, the horizontal length, being the $4$ in the $4 : 3$ ratio, is simply $\frac{4}{5}$ times the hypotenuse. $\frac{4}{5}\cdot27=21.6 \approx \boxed{\textbf{(D) } 21.5}$.

 


	\item (\textit{AMC 10 10B 2003}) The symbolism $\lfloor x \rfloor$ denotes the largest integer not exceeding $x$. For example, $\lfloor 3 \rfloor = 3,$ and $\lfloor 9/2 \rfloor = 4$. Compute

 \[\lfloor \sqrt{1} \rfloor + \lfloor \sqrt{2} \rfloor + \lfloor \sqrt{3} \rfloor + \cdots + \lfloor \sqrt{16} \rfloor.\]


 $\textbf{(A) } 35 \qquad\textbf{(B) } 38 \qquad\textbf{(C) } 40 \qquad\textbf{(D) } 42 \qquad\textbf{(E) } 136$


 
 \textbf{Answer: }B



\textbf{Solution 1}

The first three values in the sum are equal to $1,$ the next five equal to $2,$ the next seven equal to $3,$ and the last one equal to $4.$ For example, since $2^2=4$ any square root of a number less than $4$ must be less than $2.$ Sum them all together to get

 
 \[3\cdot1 + 5\cdot2 + 7\cdot3 + 1\cdot4 = 3+10+21+4 = \boxed{\textbf{(B) \ } 38}\]

 


 



\textbf{Solution 2}

We are asked to compute: 
 \[\sum_{k=1}^{16} \lfloor \sqrt{k} \rfloor\]

 ---

 
                \begin{itemize}
                    <li>
                \begin{itemize}
                    \item Step 1: Evaluate each term individually**
                \end{itemize}


                </li>
                \end{itemize}


                
 \[\lfloor \sqrt{1} \rfloor = 1 \\ \lfloor \sqrt{2} \rfloor = 1 \\ \lfloor \sqrt{3} \rfloor = 1 \\ \lfloor \sqrt{4} \rfloor = 2 \\ \lfloor \sqrt{5} \rfloor = 2 \\ \lfloor \sqrt{6} \rfloor = 2 \\ \lfloor \sqrt{7} \rfloor = 2 \\ \lfloor \sqrt{8} \rfloor = 2 \\ \lfloor \sqrt{9} \rfloor = 3 \\ \lfloor \sqrt{10} \rfloor = 3 \\ \lfloor \sqrt{11} \rfloor = 3 \\ \lfloor \sqrt{12} \rfloor = 3 \\ \lfloor \sqrt{13} \rfloor = 3 \\ \lfloor \sqrt{14} \rfloor = 3 \\ \lfloor \sqrt{15} \rfloor = 3 \\ \lfloor \sqrt{16} \rfloor = 4\]

 ---

 
                \begin{itemize}
                    <li>
                \begin{itemize}
                    \item Step 2: Group repeated values**
                \end{itemize}


                </li>
                \end{itemize}


                1+1+1 = 3  

 2+2+2+2+2 = 10  

 3+3+3+3+3+3+3 = 21  

 4 = 4

 ---

 
                \begin{itemize}
                    <li>
                \begin{itemize}
                    \item Step 3: Add them together**
                \end{itemize}


                </li>
                \end{itemize}


                3 + 10 + 21 + 4 = 38

 ---

 
                \begin{itemize}
                    <li>
                \begin{itemize}
                    \item Answer:**
                \end{itemize}


                </li>
                \end{itemize}


                
 \[\boxed{38}\]

 


	\item (\textit{AMC 10 10B 2003}) The second and fourth terms of a geometric sequence are $2$ and $6$. Which of the following is a possible first term?


 $\textbf{(A) } -\sqrt{3}  \qquad\textbf{(B) } -\frac{2\sqrt{3}}{3} \qquad\textbf{(C) } -\frac{\sqrt{3}}{3} \qquad\textbf{(D) } \sqrt{3} \qquad\textbf{(E) } 3$


 
 \textbf{Answer: }B



\textbf{Solution}

Let the first term be $a$ and the common ratio be $r$. Therefore, 

 
 \[ar=2\ \ (1) \qquad \text{and} \qquad ar^3=6\ \ (2)\]

 Dividing $(2)$ by $(1)$ eliminates the $a$, yielding $r^2=3$, so $r=\pm\sqrt{3}$.

 Now, since $ar=2$, $a=\frac{2}{r}$, so $a=\frac{2}{\pm\sqrt{3}}=\pm\frac{2\sqrt{3}}{3}$.

 We therefore see that $\boxed{\textbf{(B)}\ -\frac{2\sqrt{3}}{3}}$ is a possible first term.

 


	\item (\textit{AMC 10 10B 2003}) Find the value of $x$ that satisfies the equation

 \[25^{-2} = \frac{5^{48/x}}{5^{26/x} \cdot 25^{17/x}}.\]


 $\textbf{(A) } 2 \qquad\textbf{(B) } 3 \qquad\textbf{(C) } 5 \qquad\textbf{(D) } 6 \qquad\textbf{(E) } 9$


 
 \textbf{Answer: }B



\textbf{Solution}

Manipulate the powers of $5$ in order to get a clean expression.

 
 \[\frac{5^{\frac{48}{x}}}{5^{\frac{26}{x}} \cdot 25^{\frac{17}{x}}} = \frac{5^{\frac{48}{x}}}{5^{\frac{26}{x}} \cdot 5^{\frac{34}{x}}} = 5^{\frac{48}{x}-(\frac{26}{x}+\frac{34}{x})} = 5^{-\frac{12}{x}}\]
 \[25^{-2} = (5^2)^{-2} = 5^{-4}\]
 \[5^{-4} = 5^{-\frac{12}{x}}\]

 If two numbers are equal, and their bases are equal, then their exponents are equal as well.  Set the two exponents equal to each other.

 
 \begin{align*}-4&=\frac{-12}{x}\\ -4x&=-12\\ x&=\boxed{\textbf{(B) \ } 3}\end{align*}


 


	\item (\textit{AMC 10 10B 2003}) Nebraska, the home of the AMC, changed its license plate scheme. Each old license plate consisted of a letter followed by four digits. Each new license plate consists of three letters followed by three digits. By how many times is the number of possible license plates increased?


 $\textbf{(A) } \frac{26}{10} \qquad\textbf{(B) } \frac{26^2}{10^2} \qquad\textbf{(C) } \frac{26^2}{10} \qquad\textbf{(D) } \frac{26^3}{10^3} \qquad\textbf{(E) } \frac{26^3}{10^2}$


 
 \textbf{Answer: }C



\textbf{Solution}

There are $26$ letters and $10$ digits. There were $26 \cdot 10^4$ old license plates. There are $26^3 \cdot 10^3$ new license plates. The number of license plates increased by

 
 \[\frac{26^3 \cdot 10^3}{26 \cdot 10^4} = \boxed{\textbf{(C) \ } \frac{26^2}{10}}\]

 Note: Nebraska is not actually the home of the AMC, the AMC resides in Washington D.C., where its headquarters is located.

 


	\item (\textit{AMC 10 10B 2003}) A line with slope $3$ intersects a line with slope $5$ at point  $(10,15)$. What is the distance between the $x$-intercepts of these two lines?


 $\textbf{(A) } 2 \qquad\textbf{(B) } 5 \qquad\textbf{(C) } 7 \qquad\textbf{(D) } 12 \qquad\textbf{(E) } 20$


 
 \textbf{Answer: }A



\textbf{Solution 1}

Using the point-slope form, the equation of each line is

 
 \[y-15=3(x-10) \longrightarrow y=3x-15\]
 \[y-15=5(x-10) \longrightarrow y=5x-35\]

 Substitute in $y=0$ to find the $x$-intercepts.

 
 \[0=3x-15\longrightarrow x=5\]
 \[0=5x-35\longrightarrow x=7\]The difference between them is $7-5=\boxed{\textbf{(A) \ } 2}$.

 



\textbf{Solution 2(using slopes)}

The $x$-intercepts of a line is where the line's $y$-coordinate is $=0$. The slope is defined as $\frac{\text{change in } y\text{-coordinate}}{\text{change in } x\text{-coordinate}}$. Therefore, the line with slope $3$ can be expressed as

 
 \[\frac{\text{change in }y\text{-coordinate}}{\text{change in }x\text{-coordinate}}=3\]

 From $(10,15)$ to the $x$-intercept, the line must move $-15$ units to the $x$-axis. Therefore, 
 \[\frac{\text{change in }y\text{-coordinate}}{\text{change in }x\text{-coordinate}}=\frac{-15}{x_1}=3\implies x_1=-5\].

 Therefore, the $x$-intercept of the line with slope $3$ is $(10-5,0)=(5,0)$.Note that some sources may state "the $x$-intercept is $5$" instead of "the $x$-intercept is $(5,0)$.

 If an idea functions great for one part of the problem but does not find the solution to the problem, we want to use the idea again. Using the idea for the line with slope $5$, we find
 \[\frac{\text{change in }y\text{-coordinate}}{\text{change in }x\text{-coordinate}}=\frac{-15}{x_2}=5\implies x_2=-3\].Therefore, the $x$-intercept of the line with slope $5$ is $(10-3,0)=(7,0)$. The distance between $\boxed{(5,0)--(7,0)=2\implies A}$.

 


	\item (\textit{AMC 10 10B 2003}) Al, Betty, and Clare split $\$$1000$ among them to be invested in different ways. Each begins with a different amount. At the end of one year, they have a total of $\$$1500$. Betty and Clare have both doubled their money, whereas Al has managed to lose $\$$100$. What was Al's original portion?


 $\textbf{(A) }\$$250 \qquad\textbf{(B) }\$ $350 \qquad\textbf{(C) }\$$400\qquad\textbf{(D) }\$ $450\qquad\textbf{(E) }\$$500$


 
 \textbf{Answer: }C



\textbf{Solution 1}

For this problem, we will have to write a three-variable equation, but not necessarily solve it. Let $a, b,$ and $c$ represent the original portions of Al, Betty, and Clare, respectively. At the end of one year, they each have $a-100, 2b,$ and $2c$. From this, we can write two equations. 

 
 \[a+b+c=1000 \implies 2a+2b+2c=2000\]
 \[(a-100)+2b+2c=1500 \implies a+2b+2c=1600\]

 Since all we need to find is $a,$ subtract the second equation from the first equation to get $a=400.$ Al's original portion was $\boxed{\textbf{(C)}\ \textdollar 400}$.

 



\textbf{Solution 2}

Suppose Betty and Clare $1000-x$ dollars and Al $x$ dollars originally. Then, $(x-100)+2(1000-x)=1500 \implies x=400$, so Al's original portion was $\boxed{\textbf{(C)}\ \textdollar 400}$.

 



\textbf{Solution 3}

If Al had not lost $\textdollar 100$, the total amount would be $\textdollar 1600$. So, Betty and Clare collectively gained $\textdollar 600$ and as they have doubled their collective fortune,they must have $\textdollar 600$ with them at the beginning. This leaves $\boxed{\textbf{(C)}\ \textdollar 400}$ for Al.~Anshulb

 



\textbf{Solution 4}

Let $a, b,$ and $c$ represent the original portions of Al, Betty, and Clare, respectively. We can write some equations:
 \[a+b+c = 1000\]
 \[2b+2c+a-100 = 1500 \Longrightarrow 2b+2c+a=1600\]Subtracting the first equation from the second we get 
 \[b+c = 600\]. So, 
 \[a = 1000-600 = \boxed{\textbf{(C)}\ \textdollar 400}\]

 -jb2015007

 


	\item (\textit{AMC 10 10B 2003}) Let $\clubsuit(x)$ denote the sum of the digits of the positive integer $x$. For example, $\clubsuit(8)=8$ and $\clubsuit(123)=1+2+3=6$. For how many two-digit values of $x$ is $\clubsuit(\clubsuit(x))=3$?


 $\textbf{(A) } 3 \qquad\textbf{(B) } 4 \qquad\textbf{(C) } 6 \qquad\textbf{(D) } 9 \qquad\textbf{(E) } 10$


 
 \textbf{Answer: }E



\textbf{Solution}

Let $y=\clubsuit (x)$. Since $x \leq 99$, we have $y \leq 18$. Thus if $\clubsuit (y)=3$, then $y=3$ or $y=12$. The 3 values of $x$ for which $\clubsuit (x)=3$ are 12, 21, and 30, and the 7 values of $x$ for which $\clubsuit (x)=12$ are 39, 48, 57, 66, 75, 84, and 93. There are $\boxed{E}$ = $\boxed{10}$ values in all.

 


	\item (\textit{AMC 10 10B 2003}) Given that $3^8\cdot5^2=a^b,$ where both $a$ and $b$ are positive integers, find the smallest possible value for $a+b$.


 $\textbf{(A) } 25 \qquad\textbf{(B) } 34 \qquad\textbf{(C) } 351 \qquad\textbf{(D) } 407 \qquad\textbf{(E) } 900$


 
 \textbf{Answer: }D



\textbf{Solution}


 \[3^8\cdot5^2 = (3^4)^2\cdot5^2 = (3^4\cdot5)^2 = 405^2\]

 $405$ is not a perfect power, so the smallest possible value of $a+b$ is $405+2=\boxed{\textbf{(D)}\ 407}$.

 


	\item (\textit{AMC 10 10B 2003}) There are $100$ players in a single tennis tournament. The tournament is single elimination, meaning that a player who loses a match is eliminated. In the first round, the strongest $28$ players are given a bye, and the remaining $72$ players are paired off to play. After each round, the remaining players play in the next round. The tournament continues until only one player remains unbeaten. The total number of matches played is


 $\textbf{(A) } \text{a prime number} \qquad\textbf{(B) } \text{divisible by 2} \qquad\textbf{(C) } \text{divisible by 5} \qquad\textbf{(D) } \text{divisible by 7} \qquad\textbf{(E) } \text{divisible by 11}$


 
 \textbf{Answer: }E



\textbf{Solution 1}

Notice that $99$ players need to be eliminated for there to be declared a winner. Notice also that every match eliminates exactly one person. Therefore, $99$ matches are needed to eliminate $99$ people and therefore declare a winner. The rest of the information is irrelevant. Therefore, the total number of matches is $\text{divisible by 11}, \boxed{\text{E}}$.

 



\textbf{Solution 2}

28 people receive byes, so in the first round there are $36$ matches played. In the second round there are $28 + 36 = 64$ people, so there are 32 matches. In the subsequent rounds, there are $16, 8, 4, 2, 1$ matches played, for a total of $36 + 32 + 16 + 8 + 4 + 2 + 1 = 99$ matches. Therefore, the answer is $\rightarrow\textbf{(E) }\text{divisible by 11.}$

 


	\item (\textit{AMC 10 10B 2003}) A restaurant offers three desserts, and exactly twice as many appetizers as main courses. A dinner consists of an appetizer, a main course, and a dessert. What is the least number of main courses that a restaurant should offer so that a customer could have a different dinner each night in the year $2003$?


 $\textbf{(A) } 4 \qquad\textbf{(B) } 5 \qquad\textbf{(C) } 6 \qquad\textbf{(D) } 7 \qquad\textbf{(E) } 8$


 
 \textbf{Answer: }E



\textbf{Solution 1}

Let $m$ be the number of main courses the restaurant serves, so $2m$ is the number of appetizers. Then the number of dinner combinations is $2m\times m\times3=6m^2$. Since the customer wants to eat a different dinner in all $365$ days of $2003$, we must have

 
 \begin{align*} 6m^2 &\geq 365\\ m^2 &\geq 60.83\ldots.\end{align*}
Also, year 2003 is not a leap year, because 2003 divided by 4 does not equal an integer.The smallest integer value that satisfies this is $\boxed{\textbf{(E)}\ 8}$.

 



\textbf{Solution 2}

Let $m$ denote the number of main courses needed to meet the requirement. Then the number of dinners available is $3\cdot m \cdot 2m = 6m^2$. Thus $m^2$ must be at least $365/6 \approx 61$. Since $7^2 = 49<61<64 = 8^2$, $\boxed{8}$ main courses is enough, but 7 is not. The smallest integer value that satisfies this is $\boxed{\textbf{(E)}\ 8}$.

 


 



\textbf{Solution 3}

Let $x$ be the number of main courses and let $2x$ be the number of appetizers.Since there are 3 desserts, the number of possible dinner choices would be $2x \cdot x \cdot 3 = 6x^2$ for any number $x$. Since a year has $365$ days, we can assume that:
 \begin{align*} 6x^2 \ge 365 \end{align*}

 \begin{align*} x^2 \ge 61 \end{align*}

 \begin{align*} x \ge 7.8 \end{align*}
The least option that is greater than $7.8$ is $8$, so the answer is $\boxed{\textbf{(E)}\ 8}$.

 


	\item (\textit{AMC 10 10B 2003}) An ice cream cone consists of a sphere of vanilla ice cream and a right circular cone that has the same diameter as the sphere. If the ice cream melts, it will exactly fill the cone. Assume that the melted ice cream occupies $75\%$ of the volume of the frozen ice cream. What is the ratio of the cone's height to its radius? (Note: a cone with radius $r$ and height $h$ has volume $\pi r^2 h / 3$ and a sphere with radius $r$ has volume $4 \pi r^3 / 3$.)


 $\textbf{(A) } 2:1 \qquad\textbf{(B) } 3:1 \qquad\textbf{(C) } 4:1 \qquad\textbf{(D) } 16:3 \qquad\textbf{(E) } 6:1$


 
 \textbf{Answer: }B



\textbf{Solution}

Let $r$ be the radius of both the cone and the sphere of ice cream, and let $h$ be the height of the cone. The cone can hold a maximum of $\pi r^2h/3$ melted ice cream. The volume of the frozen ice cream is $4\pi r^3/3.$ The volume of the melted ice cream is $75 \%$ of that, which is $\pi r^3.$ Since the melted ice cream fills the cone exactly,

 
 \begin{align*}\pi r^3&=\pi r^2h/3\\ r&=h/3\\ 3&=h/r\end{align*}


 Then the ratio of the cone's height to its radius is $h:r=\boxed{\textbf{(B)}\ 3:1}$

 


	\item (\textit{AMC 10 10B 2003}) What is the largest integer that is a divisor of

 \[(n+1)(n+3)(n+5)(n+7)(n+9)\]
for all positive even integers $n$?


 $\textbf{(A) } 3 \qquad\textbf{(B) } 5 \qquad\textbf{(C) } 11 \qquad\textbf{(D) } 15 \qquad\textbf{(E) } 165$


 
 \textbf{Answer: }D



\textbf{Solution 1}

For all consecutive odd integers, one of every five is a multiple of 5 and one of every three is a multiple of 3. The answer is $3 \cdot 5 = 15$, so  ${\boxed{\textbf{(D)15}}}$ is the correct answer.

 



\textbf{Solution 2}

We'll just test all the answer choices.

 Note that for any 3 consecutive odd integers, there will be exactly one multiple of $3.$

 Let's list all possibilities of 3 consecutive odd integers. (multiple of 3, not multiple of 3, not multiple of 3), (not multiple of 3, multiple of 3, not multiple of 3) and (not multiple of 3, not multiple of 3, multiple of 3)

 To support this further, list the first few  consecutive lists of 3 consecutive odd integers.

 We have $(1, 3, 5), (3, 5, 7), (5, 7, 9), (7, 9, 11), (11, 13, 15), \ldots$

 So if the first two consecutive odd numbers aren't multiples of three, the last one must be a multiple of three.

 Therefore for five consecutive odd integers, there must be at least one multiple of three.

 In the same fashion as we did above, note that for any 5 consecutive integers, there will also be exactly one multiple of $5.$

 Therefore, for any 5 consecutive odd integers, there must be exactly one multiple of five.

 We can skip 7 since none of the answer choices are a multiple of 7.

 Now we try $11.$ 11 doesn't work since as we see the following set of 5 consecutive odd integers doesn't fit: $(13, 15, 17, 19, 21).$

 Since any 5 consecutive integers is divisible both by $3$ and $5$, it also must be divisible by ${\boxed{\textbf{(D)15}}}$ and no higher since we saw that $11$ does not work and that there is no answer choice that is multiple of $7$.

 


 



\textbf{Solution 3}

For these types of questions, I like to test out two to three values of $n$ and look at the prime factorizations. I tried $n = 2$ and $n = 14$ so that there would be no overlapping of numbers. When $n = 2$, you get $3 * 5 * 7 * 11$ which has a prime factorization of $3^3 * 5 * 7 * 9 * 11$. When $n = 14$, you get $15 * 17 * 19 * 21 * 23$. Because the problem states it must be valid for all positive even integers, whichever factor is shared between the two cases will be the answer. The prime factorization of $n = 14$ is $5 * 3^2 * 17 * 19 * 7 * 23$. Notice that the common factor is $5*3$, or $15$. So, our answer is ${\boxed{\textbf{(D)15}}}$.

 -N.

 



\textbf{Solution 4 (Mods, more rigorous)}

The answer choices give us three possible divisors to investigate: 3, 5, and 11; 15 and 165 are products of some of those three. Observe that 2, the common difference between the five factors, is coprime with 3, 5, and 11. This means that for any given residue for $n+1$, as we have five factors, we are guaranteed that at least one of the residues modulo 3 and modulo 5 is 0, so at least one factor has residue 0 modulo 3 and modulo 5. (We can confirm this by a little casework, or we can see that even with a $n+1$ as a 'generator', it still will cycle within five times of adding 2). This guarantees 3 and 5 as factors. Yet, there are too many residues in modulo 11 for these factors to cycle through, so none of the $n+1,n+3,\dots,n+9$ is guaranteed to have a factor of 11. Thus, ${\boxed{\textbf{(D)15}}}$.

 In fact, using this reasoning, we can show that any prime greater than 5 isn't a guaranteed factor. Furthermore, since these integers are all odd (alternatively, because the difference is 2, we don't run through all the residues modulo 2 or 4), none of them have factor 2, so we are sure that 15 is the largest possible value that is guaranteed - no answer choice analysis needed!

 


	\item (\textit{AMC 10 10B 2003}) Three semicircles of radius $1$ are constructed on diameter $\overline{AB}$ of a semicircle of radius $2$. The centers of the small semicircles divide $\overline{AB}$ into four line segments of equal length, as shown. What is the area of the shaded region that lies within the large semicircle but outside the smaller semicircles?


 
                    \begin{figure}[H]
                        \centering
                        \includegraphics[width=0.5\textwidth]{Problem/AMC_10_10B_2003_19_Problem_Figure_1.png}
                    \end{figure}
                    


 $\textbf{(A) } \pi - \sqrt{3} \qquad\textbf{(B) } \pi - \sqrt{2} \qquad\textbf{(C) } \frac{\pi + \sqrt{2}}{2} \qquad\textbf{(D) } \frac{\pi +\sqrt{3}}{2} \qquad\textbf{(E) } \frac{7}{6}\pi - \frac{\sqrt{3}}{2}$


 
 \textbf{Answer: }E



\textbf{Solution 1}


                    \begin{figure}[H]
                        \centering
                        \includegraphics[width=0.5\textwidth]{Solution/AMC_10_10B_2003_19_Solution_Figure_1.png}
                    \end{figure}
                    

 By drawing four lines from the intersect of the semicircles to their centers, we have split the white region into $\frac{5}{6}$ of a circle with radius $1$ and two equilateral triangles with side length $1$. This gives the area of the white region as $\frac{5}{6}\pi+\frac{2\cdot\sqrt3}{4}=\frac{5}{6}\pi+\frac{\sqrt3}{2}$. The area of the shaded region is the area of the white region subtracted from the area of the large semicircle. This is equivalent to $2\pi-\left(\frac{5}{6}\pi+\frac{\sqrt3}{2}\right)=\frac{7}{6}\pi-\frac{\sqrt3}{2}$.

 Thus the answer is $\boxed{\textbf{(E)}\ \frac{7}{6}\pi-\frac{\sqrt3}{2}}$.

 \textbf{Note} The reason why it is $\frac{5}{6}$ of a circle and why the triangles are equilateral are because, first, the radii are the same and they make up the equilateral triangles.

 Secondly, the reason it is $\frac{5}{6}$ of a circle is because the middle sector has a degree of $180-2 \cdot 60 = 60$ and thus $\frac{60}{360}=\frac{1}{6}$ of a circle.

 The other two have areas of $\frac{180-60}{360}=\frac{1}{3}$ of a triangle each.

 Therefore, the total fraction of the circle(since they have the same radii) is $\frac{1}{6} + 2 \cdot \frac{1}{3} = \frac{1}{6} + \frac{4}{6} = \frac{5}{6}.$

 



\textbf{Solution 2(Answer Choices)}

Answer choices A and B are impossible since the area is obviously not a circle of radius 1 minus a triangle.

 Answer choices C and D are also impossible since they are adding but we are subtracting.

 That leaves us with only answer choice $\textbf{(E)}.$

 



\textbf{Solution 3 ( Circle area formulas )}

We can see that the area of the $3$ overlapped white semi circles is equal to :

 $2*\frac{1}{2}\pi+\frac{1}{2}*\frac{\pi}{3}-2*\frac{1}{2}\left(\frac{\pi}{3} -\ \frac{\sqrt{3}}{2}\right)$

 We get this because we think of the 2 overlapped circles on the side as 2 full semi circles then adding the area of the sector in between but subtracting the area segments of the 2 semi circles out (So that they don't overlap). Now by using the formulas:

 Area of sector $=\frac{1}{2}r^2*\text{sin}\theta$

 Area of segment $=\frac{1}{2}r^2\left( \theta -\ \text{sin}\theta\right)$

 ($\theta$ is in radians)

 and so we get :

 $=\frac{5}{6}\pi+\frac{\sqrt3}{2}$

 then the solution is the same as in solution 1.

 


	\item (\textit{AMC 10 10B 2003}) In rectangle $ABCD, AB=5$ and $BC=3$. Points $F$ and $G$ are on $\overline{CD}$ so that $DF=1$ and $GC=2$. Lines $AF$ and $BG$ intersect at $E$. Find the area of $\triangle AEB$.


 
                    \begin{figure}[H]
                        \centering
                        \includegraphics[width=0.5\textwidth]{Problem/AMC_10_10B_2003_20_Problem_Figure_1.png}
                    \end{figure}
                    
$\textbf{(A) } 10 \qquad\textbf{(B) } \frac{21}{2} \qquad\textbf{(C) } 12 \qquad\textbf{(D) } \frac{25}{2} \qquad\textbf{(E) } 15$


 
 \textbf{Answer: }D



\textbf{Solution 1}

$\triangle EFG \sim \triangle EAB$ because $FG \parallel AB.$ The ratio of $\triangle EFG$ to $\triangle EAB$ is $2:5$ since $AB=5$ and $FG=2$ from subtraction. If we let $h$ be the height of $\triangle EAB,$

 
 \[\frac{2}{5} = \frac{h-3}{h}\]
 \[2h = 5h-15\]
 \[3h-15 = 0\]
 \[3h = 15\]
 \[h = 5\]

 The height is $5$ so the area of $\triangle EAB$ is $\frac{1}{2}(5)(5) = \boxed{\textbf{(D)}\ \frac{25}{2}}$.

 



\textbf{Solution 2}

We can look at this diagram as if it were a coordinate plane with point $A$ being $(0,0)$. This means that the equation of the line $AE$ is $y=3x$ and the equation of the line $EB$ is $y=\frac{-3}{2}x+\frac{15}{2}$. From this we can set of the follow equation to find the $x$ coordinate of point $E$:

 
 \[3x=\frac{-3}{2}x+\frac{15}{2}\]
 \[6x=-3x+15\]
 \[9x=15\]
 \[x=\frac{5}{3}\]

 We can plug this into one of our original equations to find that the $y$ coordinate is $5$, meaning the area of $\triangle EAB$ is $\frac{1}{2}(5)(5) = \boxed{\textbf{(D)}\ \frac{25}{2}}$

 



\textbf{Solution 3}

At points $A$ and $B$, segment $AE$ is 5 units from segment $BE$. At points $F$ and $G$, the segments are 2 units from each other. This means that collectively, the two lines closed the distance between them by 3 units over a height of 3 units. Therefore, to close the next two units of distance, they will have to travel a height of 2 units.

 Then calculate the area of trapezoid $FGBA$ and triangle $EGF$ separately and add them. The area of the trapezoid is $\frac {2+5}{2}\cdot 3 = \frac {21}{2}$ and the area of the triangle is $\frac{1}{2}\cdot 2 \cdot 2 = 2$. $\frac{21}{2}+2=\boxed{\textbf{(D)}\ \frac{25}{2}}$

 



\textbf{Solution 4}

Since $\Delta{ABE}\sim{\Delta{FGE}}$ then $[AFGB]\sim{[FXYG]}$, where $X$ and $Y$ are ponts on $EF$ and $EG$ respectivley which make the areas similar. This process can be done over and over again multiple times by the ratio of $\frac{FG}{AB}=\frac{2}{5}$, or something like this
 \[[AEB]=[AFGB]+[FXYZ]+...\]
 \[[AEB]=[AFGB]+\frac{2}{5}[AFGB]+(\frac{2}{5})^2[AFGB]+...\]we have to find the ratio of the areas when the sides have shrunk by length $\frac{2}{5}l$

 
                    \begin{figure}[H]
                        \centering
                        \includegraphics[width=0.5\textwidth]{Solution/AMC_10_10B_2003_20_Solution_Figure_1.png}
                    \end{figure}
                    

 Let $[AFGB]'$ be the area of the shape whose length is $\frac{2}{5}l$
 \[[AFGB]'=[ADCB]-[ADF]-[BCG]\]
 \[[AFGB]'=12/5-6/25-12/25\]
 \[[AFGB]'=42/25\]Now comparing the ratios of $[AFGB]'$ to $[AFGB]$ we get
 \[\frac{[AFGB]'}{[AFGB]}=\frac{42}{25}/\frac{21}{2}\implies \frac{[AFGB]'}{[AFGB]}=\frac{4}{25}\]By applying an infinite summation
 \[[AEB]=\sum_{n=0}^{\infty} \frac{21}{2}\cdot{(\frac{4}{25})^n}\]
 \[S=\frac{a_1}{1-r}\]
 \[\boxed{[AEB]=\frac{25}{2}}\]

 



\textbf{Solution 5}

Drop a perpendicular from $E$ to $AB$ and call the intersection point $X.$ $\triangle EXA$ and $\triangle DFA$ are similar and so are $\triangle BCG$ and $\triangle EXB$ (You can prove this with $AA$ by observing that $\angle DFA$ is congruent to $\angle EAB$ and so on). This means that $\dfrac{AX}{BX}=\dfrac{1}{2}$ so $AX=\dfrac{5}{3}.$ $\dfrac{EX}{AX}=3,$ which means that $EX=5.$ Then, $[AEB]=\dfrac{5\cdot 5}{2}=\boxed{\textbf{(D) } \dfrac{25}{2}}.$

 



\textbf{Solution 6}

We see that triangles $EFG$ and $EAB$ are similar by Angle-Angle, and their scale factor must be $2:5$. The area of quadrilateral $AFGB$ is simply $15 - 4.5 = 10.5$ (the area of the rectangle minus the two smaller triangles). 

 We can let the area of triangle $EFG$ equal $x$, which means the area of triangle $EAB$ is $x + 10.5$. Since we know the scale factor for the linear measures of the triangles is $2:5$, the ratio of their areas must be $k^2$. In other words, the ratio of the area of $\triangle EFG$ to the area of $\triangle EAB$ is $2^2: 5^2$, or $4:25$. 

 Now we can set up an equation to solve for $x$:
 \[\frac{25}{4}x = x + 10.5\]

 Subtracting $x$ from both sides gives:
 \[\frac{21}{4}x = 10.5\]

 Solving this equation, we get $x = 2$, meaning the area of $\triangle EFG$ is $2$. This means that the area of triangle $EAB$ is $10.5 + 2 = 12.5$, which corresponds to answer choice $\text{(D)}\ \boxed{\frac{25}{2}}$.

 -jpudi

 


	\item (\textit{AMC 10 10B 2003}) A bag contains two red beads and two green beads. You reach into the bag and pull out a bead, replacing it with a red bead regardless of the color you pulled out. What is the probability that all beads in the bag are red after three such replacements?


 $\textbf{(A) } \frac{1}{8} \qquad\textbf{(B) } \frac{5}{32} \qquad\textbf{(C) } \frac{9}{32} \qquad\textbf{(D) } \frac{3}{8} \qquad\textbf{(E) } \frac{7}{16}$


 
 \textbf{Answer: }C



\textbf{Solution 1 (Casework)}

We can divide the case of all beads in the bag being red after three replacements into three cases. 

 The first case is in which the two green beads are the first two beads to be chosen. The probability for this is
 \[\frac{2}{4} \times \frac{1}{4} \times \frac{4}{4} = \frac{1}{8}\] 

 The second case is in which the green beads are chosen first and third. The probability for this is
 \[\frac{2}{4} \times \frac{3}{4} \times \frac{1}{4} = \frac{3}{32}\]

 The third case is in which the green beads are chosen second and third. The probability for this is
 \[\frac{2}{4} \times \frac{2}{4} \times \frac{1}{4} = \frac{1}{16}\]

 Add all these cases together
 \[\frac{1}{8}+\frac{3}{32}+\frac{1}{16} = \frac{4}{32}+\frac{3}{32}+\frac{2}{32} = \boxed{\textbf{(C) \ } \frac{9}{32}}\]

 



\textbf{Solution 2 (Probability States)}

We use probability states. Notice that we can construct a Markov chain like the one below.

 
                    \begin{figure}[H]
                        \centering
                        \includegraphics[width=0.5\textwidth]{Solution/AMC_10_10B_2003_21_Solution_Figure_1.png}
                    \end{figure}
                    From this, each arrow represents one distinct replacement. Then, one can notice that we can loop from 2 Red to 2 Red, and then move from 2 Red to 3 Red and from 3 Red to All red. However another case is to go from 2 Red to 3 Red directly, and loop once at 3 Red, and move to 4 Red. The third case is to immediately in two moves go from 2 Red to 3 Red, and then from 3 Red to 4 Red, and to loop at 4 Red.

 The first case has a probability $\frac{1}{2} \cdot \frac{1}{2} \cdot \frac{1}{4} = \frac{1}{16}$. The second case has a probability $\frac{1}{2} \cdot \frac{3}{4} \cdot \frac{1}{4} = \frac{3}{32}$. The third case has probability $\frac{1}{2} \cdot \frac{1}{4} \cdot 1 = \frac{1}{8}$.

 The asked probability is $\frac{2}{32} + \frac{3}{32} + \frac{4}{32} = \boxed{\frac{9}{32}}$.

 


	\item (\textit{AMC 10 10B 2003}) A clock chimes once at $30$ minutes past each hour and chimes on the hour according to the hour. For example, at $1 \text{PM}$ there is one chime and at noon and midnight there are twelve chimes. Starting at $11:15 \text{AM}$ on $\text{February 26, 2003},$ on what date will the $2003\text{rd}$ chime occur?


 $\textbf{(A) } \text{March 8} \qquad\textbf{(B) } \text{March 9} \qquad\textbf{(C) } \text{March 10} \qquad\textbf{(D) } \text{March 20} \qquad\textbf{(E) } \text{March 21}$


 
 \textbf{Answer: }B



\textbf{Solution}

First, find how many chimes will have already happened before midnight (the beginning of the day) of $\text{February 27, 2003}.$ $13$ half-hours have passed, and the number of chimes according to the hour is $1+2+3+\cdots+12.$ The total number of chimes is $13+78=91.$

 Every day, there will be $24$ half-hours and $2(1+2+3+\cdots+12)$ chimes according to the arrow, resulting in $24+156=180$ total chimes.

 On $\text{February 26},$ the number of chimes that still need to occur is $2003-91=1912.$ $1912 \div 180=10 \text{R}112.$ Rounding up, it is $11$ days past $\text{February 26},$ which is $\boxed{\textbf{(B) \ } \text{March 9}}$

 


	\item (\textit{AMC 10 10B 2003}) A regular octagon $ABCDEFGH$ has an area of one square unit. What is the area of the rectangle $ABEF$?


 
                    \begin{figure}[H]
                        \centering
                        \includegraphics[width=0.5\textwidth]{Problem/AMC_10_10B_2003_23_Problem_Figure_1.png}
                    \end{figure}
                    
$\textbf{(A) } 1 - \frac{\sqrt{2}}{2} \qquad\textbf{(B) } \frac{\sqrt{2}}{4} \qquad\textbf{(C) } \sqrt{2} - 1 \qquad\textbf{(D) } \frac{1}{2} \qquad\textbf{(E) } \frac{1+\sqrt{2}}{4}$


 
 \textbf{Answer: }D



\textbf{Solution 1}

Here is an easy way to look at this, where $p$ is the perimeter, and $a$ is the \href{https://artofproblemsolving.com/wiki/index.php?title=Apothem}{apothem}:

 Area of Octagon: $\frac{ap}{2}=1$.

 Area of Rectangle: $\frac{p}{8}\times 2a=\dfrac{2ap}{8}=\frac{ap}{4}$.

 You can see from this that the octagon's area is twice as large as the rectangle's area is $\boxed{\textbf{(D)}\ \frac{1}{2}}$.

 



\textbf{Solution 2}


                    \begin{figure}[H]
                        \centering
                        \includegraphics[width=0.5\textwidth]{Solution/AMC_10_10B_2003_23_Solution_Figure_1.png}
                    \end{figure}
                    

 Here is a less complicated way than that of the user below. If you draw a line segment from each vertex to the center of the octagon and draw the rectangle ABEF (in red), you can see that $2$ of the triangles (in blue) share the same base and height with $\dfrac{1}{2}$ the rectangle. Therefore, the rectangle's area is the same as $2\cdot2$ of the $8$ triangles, and is  $\boxed{\textbf{(D)}\ \frac{1}{2}}$ the area of the octagon.

 



\textbf{Solution 3}


                    \begin{figure}[H]
                        \centering
                        \includegraphics[width=0.5\textwidth]{Solution/AMC_10_10B_2003_23_Solution_Figure_2.png}
                    \end{figure}
                    Drawing lines $AD$, $BG$, $CF$, and $EH$, we can see that the octagon is comprised of $1$ square, $4$ rectangles, and $4$ triangles. The triangles each are $45-45-90$ triangles, and since their diagonal is length $x$, each of their sides is $\frac{\sqrt{2}}{2}x$. The area of the entire figure is, likewise, $x^2$ (the square)$+4x^2\frac{\sqrt{2}}{2}$ (the 4 rectangles)$+2\cdot(\frac{\sqrt{2}}{2}x)^2$ (the triangles), which simplifies to $2x^2 + 2\sqrt{2}x^2$. The area of $ABEF$ is just $x(x+\frac{2\sqrt{2}}{2}x)$, or $x^2$ + $x^2\sqrt{2}$, which we can see is the area of $\frac{ABCDEFGH}{2} =  \boxed{\textbf{(D)}\ \frac{1}{2}}$ the area of the octagon.

 



\textbf{Solution 4}

First, we are going to divide the diagram. Now we need to find the ratio of the area of the rectangle to the area of the trapezoid. 

 
                    \begin{figure}[H]
                        \centering
                        \includegraphics[width=0.5\textwidth]{Solution/AMC_10_10B_2003_23_Solution_Figure_3.png}
                    \end{figure}
                    

 The area of a trapezoid is $\frac{b_1 + b_2}{2}h$Note that the trapezoid is made out of 2 45-45-90 triangles and a rectangle, and that $AH=FG=1$. By realizing that, the area of the trapezoid is $(\frac{2+\sqrt{2}}{2})(\frac{\sqrt{2}}{2}$). To make this product easier, note there is two trapezoids, so the new product is now this, $(\frac{2+\sqrt{2}}{2})(\sqrt{2}) = \sqrt{2} + 1$ 

 Notice how the rectangle has side lengths $\sqrt{2}+1$ and $1$, so it's area is also $\sqrt{2}+1$.

 The ratio of the area of rectangle $ABEF$ to the two trapezoids is $1:1$, meaning they share half the area of the octagon. Since the area of the octagon is 1, $\therefore$ the area of the rectangle is $\boxed{\textbf{(D)}\ \frac{1}{2}}$.

 


	\item (\textit{AMC 10 10B 2003}) The first four terms in an arithmetic sequence are $x + y, x - y, xy,$ and $x/y,$ in that order. What is the fifth term? 


 $\textbf{(A)}\ -\frac{15}{8}\qquad\textbf{(B)}\ -\frac{6}{5}\qquad\textbf{(C)}\ 0\qquad\textbf{(D)}\ \frac{27}{20}\qquad\textbf{(E)}\ \frac{123}{40}$


 
 \textbf{Answer: }E



\textbf{Solution 1}

The difference between consecutive terms is $(x-y)-(x+y)=-2y.$ Therefore we can also express the third and fourth terms as $x-3y$ and $x-5y.$ Then we can set them equal to $xy$ and $\frac{x}{y}$ because they are the same thing.

 
 \begin{align*} xy&=x-3y\\ xy-x&=-3y\\ x(y-1)&=-3y\\ x&=\frac{-3y}{y-1} \end{align*}


 Substitute into our other equation.

 
 \[\frac{x}{y}=x-5y\]
 \[\frac{-3}{y-1}=\frac{-3y}{y-1}-5y\]
 \[-3=-3y-5y(y-1)\]
 \[0=5y^2-2y-3\]
 \[0=(5y+3)(y-1)\]
 \[y=-\frac35, 1\]

 But $y$ cannot be $1$ because then the first term would be $x+1$ and the second term $x-1$ while the last two terms would be equal to $x.$ Therefore $y=-\frac35.$ Substituting the value for $y$ into any of the equations, we get $x=-\frac98.$ Finally,

 
 \[\frac{x}{y}-2y=\frac{9\cdot 5}{8\cdot 3}+\frac{6}{5}=\boxed{\textbf{(E)}\ \frac{123}{40}}\]

 



\textbf{Solution 2}

Because this is an arithmetic sequence, we conclude from the first two terms that the common difference is $-2y$. Therefore, $xy = x-3y$ and $\frac{x}{y}=x-5y$. If we multiply $xy$ and $\frac{x}{y}$, we see: 

 $xy\cdot\frac{x}{y} = (x-3y)(x-5y) = x^2 - 8xy + 15y^2$

 Because $xy\cdot\frac{x}{y}$, by basic multiplication, is $x^2$, we have 

 $x^2 = x^2 - 8xy + 15y^2$

 $8xy = 15y^2$

 $8x = 15y$

 $x = \frac{15y}{8}$

 Now that we have $x$ in terms of $y$, we substitute in $\frac{15y}{8}$ in for $x$ in $\frac{x}{y}$ (the fourth term). This leaves us with $\frac{\frac{15y}{8}}{y} = \frac{15}{8}$. 

 Recall that $\frac{x}{y}$ can be written as $x - 5y$. Thus, $x - 5y = \frac{15}{8}$. Substitute in $\frac{15}{8}$ in for $x$, and we see:

 $\frac{15y}{8} - 5y = \frac{15}{8}$

 $15y - 40y = 15$

 $y = \frac{15}{-25} = \frac{-3}{5}$

 Aha! This means $2y$, the common difference, is $2\cdot y = 2\cdot \frac{-3}{5} = \frac{-6}{5}$. Now, all we need to do is find the fifth term, which is just $\frac{x}{y}-2y$. We can substitute known values to solve: 

 $\frac{x}{y}-2y = \frac{15}{8} - \frac{-6}{5} = \frac{15}{8} + \frac{6}{5} = \frac{75 + 48}{40} = \boxed{\textbf{(E)}\ \frac{123}{40}}$.

 


 



\textbf{Solution 3}

From the first two terms, we can figure out the common difference, $x-y-(x+y)=-2y$.This means that the third, fourth and fifth terms are $x-3y$, $x-5y$ and $x-7y$ respectively.

 The third term is also equal to $xy$, so $xy=x-3y$, which can be rearranged to $xy+3y=x$.

 Dividing by y, we have$x+3=\frac{x}{y}$.

 $\frac{x}{y}$ happens to be the fourth term. Therefore,$\frac{x}{y}=x-5y=x+3$.

 $x-5y=x+3$

 $-5y=3$

 $y=-\frac{3}{5}$

 Now that we have $y$, we can substitute it into an equation above, like $\frac{x}{y}=x+3$.

 $\frac{x}{\frac{-3}{5}}=x+3$

 $-\frac{5x}{3}=x+3$

 $-\frac{5x+3x}{3}=3$

 $-\frac{8x}{3}=3$

 $8x=-9$

 $x=\frac{-9}{8}$

 As mentioned earlier, the fifth term we want in the end is equal to $x-7y$. Substitute some more and...

 $\frac{-9}{8}-7(\frac{-3}{5})=\frac{-9\cdot5}{8\cdot5}+\frac{-3\cdot-7\cdot8}{5\cdot8}=\frac{-45}{40}+\frac{168}{40}=\boxed{\textbf{(E)}\ \frac{123}{40}}$

 (Alternatively, like the above solutions suggest, you could also use the admittedly easier $\frac{x}{y}-2y$ as the final step.)

 


	\item (\textit{AMC 10 10B 2003}) How many distinct four-digit numbers are divisible by $3$ and have $23$ as their last two digits?


 $\textbf{(A) } 27 \qquad\textbf{(B) } 30 \qquad\textbf{(C) } 33 \qquad\textbf{(D) } 81 \qquad\textbf{(E) } 90$


 
 \textbf{Answer: }B



\textbf{Solution 1 (Slow)}

To test if a number is divisible by three, you add up the digits of the number. If the sum is divisible by three, then the original number is a multiple of three. If the sum is too large, you can repeat the process until you can tell whether it is a multiple of three. This is the basis for the solution. Since the last two digits are $23$, the sum of the digits is $2+3 = 5$ (therefore it is not divisible by three). However certain numbers can be added to make the sum of the digits a multiple of three.

 $5+1 = 6$,

 $5+4 = 9$, and so on. 

 However since the largest four-digit number ending with $23$ is $9923$, the maximum sum is

 $5+18 = 23$.

 Using that process we can fairly quickly compile a list of the sum of the first two digits of the number. 

 
 \[\{1, 4, 7, 10, 13, 16\}\]

 Now we find all the two-digit numbers that have any of the sums shown above. We can do this by listing all the two digit numbers $xy$ in separate cases.

 
 \[I. x+y = 1, \{10\} = 1\]
 \[II. x+y = 4, \{13, 22, 31, 40\} = 4\]
 \[III. x+y = 7, \{16, 25, 34, 43, 52, 61, 70\} = 7\]
 \[IV. x+y = 10, \{19, 28, 37, 46, 55, 64, 73, 82, 91\} = 9\]
 \[V. x+y = 13, \{49, 58, 67, 76, 85, 94\} = 6\]
 \[VI. x+y = 16, \{79, 88, 97\} = 3\]

 And finally, we add the number of elements in each set.

 
 \[1+4+7+9+6+3 = \boxed{\textbf{(B)}\ 30}\]

 



\textbf{Solution 2 (Medium)}

A number divisible by $3$ has all its digits add to a multiple of $3.$ The last two digits are $2$ and $3$ and add up to $5 \equiv 2\ (\text{mod}\ 3).$ Therefore the first two digits must add up to $1\ (\text{mod}\ 3).$ $4$ digits (including $0$) are $0\ (\text{mod}\ 3),$ $3$ are $1\ (\text{mod}\ 3),$ and $3$ are $2\ (\text{mod}\ 3).$ The following combinations are equivalent to $1\ (\text{mod}\ 3)$:

 
 \[0\ (\text{mod}\ 3)+1\ (\text{mod}\ 3) \equiv 1\ (\text{mod}\ 3)+0\ (\text{mod}\ 3) \equiv 2\ (\text{mod}\ 3) +2\ (\text{mod}\ 3)\]

 Let the first term in each combination represent the thousands digit and the second term represent the hundreds digit. We can use this to find the total amount of four-digit numbers.

 
 \[3\cdot3 + 3\cdot4 + 3\cdot3 = 9 + 12 + 9 = \boxed{\textbf{(B)}\ 30}\]

 



\textbf{Solution 3 (Fast)}

We have the following: $n \equiv 0 \pmod{3}$ and $n \equiv 23\pmod{100}$. Then $n = 3a = 100b+23$ for some integers $a$ and $b$.  Taking mod 3 gives:$0 \equiv b+2 \pmod 3 \implies b\equiv 1\pmod 3$ so $b=3c+1$ for some integer $c$.  But $N = 100b+23 = 100(3c+1)+23$ so $n = 300c+123$.  Bounding this gives us:$999 < 300c+123 < 10000$ so $876 < 300c < 9877$.  Dividing by $300$ gives $2.92 < c < 32.9222$ so $3\le c \le 32$.  This gives $32-3+1 = \boxed{\textbf{(B)} \ 30 }$

 



\textbf{Solution 4 (Extremely Fast)}

The number is in the form $xy23$

 Notice that the number is divisible by $3$ if the sum of the digits of $xy23$ is divisible by $3$

 Our first number that is divisible by $3$ is $1023$, next is $1323$..Notice $xy$ goes from $10\implies 97$

 Hence, there are $\frac{97-10}{3}+1=30$ distinct four digit numbers.

 



\textbf{Solution 5 (Even Faster than Extremely Fast)}

Following the form $xy23$ in Solution 4, notice that $x+y \equiv 1 \pmod 3$ to satisfy our condition. Choose a value of $x$ from 1 to 9. For $x=1, 4, 7$, there are exactly 4 values of $y: 0, 3, 6, 9$. For the remaining 6 digits, there are 3 choices for y. So our answer is $(3)(4)+(6)(3)=30$.

 



\textbf{Solution 6 (Even Faster than Even Faster than Extremely Fast)}

There are $9 \cdot 10 = 90$ possible values for the first two digits. $\frac{1}{3}$ of them yield a multiple of $3$, so the answer is $\frac{90}{3} = \boxed{\textbf{(B)}\ 30}$

 (In more detail: There are 90 possible values for the first two digits. All numbers can be written as 0, 1, or 2 mod(3). 2+3= 5 = -1 mod(3), so the first two digits need to add to 1 mod(3) for all the digits to equal 0 mod(3) (a multiple of 3). Out of the 3 ways to pick 2 numbers in mod(3) (0+1, 0+2, and 1+2), only $\frac{1}{3}$ of the ways (only 0+1) equals 1 mod(3). Thus, the answer is $\frac{90}{3} = \boxed{\textbf{(B)}\ 30}$.

 



\textbf{Solution 7 (Easy to understand)}

Denote the four digit number by $ab23_{10}$. Then it is well known the sum of the digits of an integer $N$ determine whether or not they are divisible by $3$. Thus we want $a + b + 2 + 3 \equiv 0 \pmod 3$.

 $\underline{a = 1}: \implies$ $b \equiv 0 \pmod 3 \implies b = 0, 3, 6, 9$.

 $\underline{a = 2}: \implies$ $b \equiv 2 \pmod 3 \implies b = 2, 5, 8$.

 $\underline{a = 3}: \implies$ $b \equiv 1 \pmod 3 \implies b = 1, 4, 7$.

 
Now we have that when $a = 4$ this is the same case as $a = 1$ as $4 \equiv 1 \pmod 3$. Hence if we count the above number of solutions and multiply by $3$ (as $1 \leq a \leq 9$), we are home free. There are $4 + 3 + 3 = 10$ solutions. Multiplying this by $3$ yields $\boxed{30}$. $\blacksquare$

 


\end{enumerate}